%%%%%%%%%%%%%%%%%%%%%%%%%%%%%%%%%%%%%%%%
% CAPÍTULO - OSIRIS-REx

\chapter{OSIRIS-REx}
\label{chp:orex}

A OSIRIS-Rex (\textit{Origins, Spectral Interpretation, Resource Identification, Security, Regolith Explorer}) é uma missão da Nasa que visa coletar uma amostra do asteroide Bennu, e trazer de volta para a Terra para ser estudado \cite{NMOR}.

\section{Missão} %e objetivo

\f{ort}{Wikipedia}{Trajetória da missão OIRIS-R}

O principal objetivo da missão é obter uma amostra de pelo menos 60 gramas do asteroide Bennu que orbita o sol e que é classificado como um objeto próximo à Terra. O material coletado será utilizado também em estudos da evolução do sistema solar, além de estudo de compostos orgânicos que levaram à formação da vida na Terra. Foi lançado em 2016 e o retorno está previsto para 2023 (\rf{ort}). O custo aproximado da missão é de aproximadamente US\$ 1 bilhão, incluindo o veículo lançador Atlas V.

A Osiris é a segunda missão a recolher amostras em asteroides, antes dela a sonda japonesa Haybusa recolheu amostras do asteroide Itokawa em 2010.
O \textit{Goddard Space Flight Center} da Nasa cuida do gerenciamento geral, engenharia, e navegação da missão. O Laboratório Lunar e Planetário da Universidade do Arizona cuida das operações científicas.
A \textit{Lockheed Martin Space Systems} construiu a espaçonave e cuida das operações da missão.
O time de cientistas conta com membros dos EUA, Canadá, França, Alemanha, e UK \cite{OSIRIS-REx}.

O lançamento ocorreu em 2016 no complexo de Cabo Canaveral, e usou um veículo lançador Atlas V. A espaçonave entrou em fase de cruzeiro assim que  separou do veículo lançador. Em seguida o painel solar foi aberto, o sistema de propulsão foi iniciado, e foi estabelecida a ligação de comunicação com a Terra. Em dezembro de 2016 realizou a sua primeira manobra no espaço profundo. Depois outra manobra em janeiro de 2017. Quando se encontrou com o asteroide em dezembro de 2018, entrou na fase de ciência e coleta de amostras.

Durante a fase de cruzeiro, ao passar próximo ao ponto L4 de Lagrange entre o Sol e a Terra, a espaçonave foi usada para estudar uma classe de objetos conhecidos como \textit{Earth-Trojan} \cite{ETrojan}. Em fevereiro de 2017 a câmera MapCam foi usada para procurar os objetos e tirar fotografias que serão estudadas pela Universidade do Arizona. Outro instrumento, a PolyCam, também foi usado para tirar fotos de Júpiter e de três luas, Calisto, Io, e Ganimedes.

Ao chegar em Bennu, a espaçonave fez várias passagens próximas à superfície do asteroide, para refinar o formato e a órbita de Bennu. Durante as passagens, o espectroscópio detectou minerais hidratados na forma de argila. Ao entrar em órbita em dezembro de 2018 a cerca de 1,75 km do asteroide, a espaçonave iniciou a campanha de mapeamento e sensoreamento remoto. No final desta etapa, a órbita passou para um raio de 1 km.

Os objetivos científicos da missão são analisar uma amostra do asteroide, mapear as propriedades dos materiais, documentar a textura, morfologia, geoquímica, e outras propriedades como força térmica, e comparação com outros asteroides.

\section{Descrição} %diagrama(espaçonave), fotos gerais, explicação geral, desenvolvimento, história da espaçonave, cronograma do projeto

\f{ord}{\citeonline{ORD}}{Principais componentes da OSIRIS-REx}
\f{ortvc}{Nasa}{OSIRIS-REx na câmara de vácuo}

As dimensões da espaçonave são aprox. 2,4 x 2,4 x 3,2 m, e pesa cerca de 2 toneladas. Com os painéis abertos mede 6,2 m. Os painéis solares geram de 1,2 a 3 kW, dependendo da distância do sol. A energia é armazenada em baterias Li-ion. A propulsão utiliza hidrazina como propelente e é baseada no sistema desenvolvido para a \textit{Mars Reconnaissance Orbiter}. Carrega 1,2 toneladas de propelente que é mantido pressurizado utilizando um tanque externo de hélio. A \rf{ortvc} mostra a OSIRIS-REx sendo colocada na câmara térmica de vácuo na Lockheed Martin para testes ambientais. 

A OSIRIS-Rex reúne todos os ingredientes necessários para realizar uma missão coleta de amostra em asteroide: a Universidade do Arizona cuida da ciência planetária e tem experiência na operação do módulo Mars Phoenix Lander; a Lockheed Martin possui experiência em desenvolvimento e operação de missões de retorno de amostra; o Goddard Space Flight Center (Nasa) possui experiência em gerenciamento de projetos, engenharia de sistemas, segurança, e garantia de missão e espectroscopia de infravermelho próximo visível; A KinetX possui experiência em navegação de espaçonaves; A Universidade do Estado do Arizona tem experiência com espectrômetros de emissão térmica; a Agência Espacial Canadense forneceu o altímetro a laser, que também foi usado na missão Phoenix Mars; o MIT e o Harvard College Observatory forneceram um espectrômetro de raios-X de imagem; além de uma equipe de cientistas dos Estados Unidos, Canadá, França, Alemanha, Grã-Bretanha e Itália
\cite{osiris2020}.

O asteroide Bennu é um objeto próximo à Terra com um diâmetro de aprox. 500 metros, massa de $ 7,8 \times 10^{10} $ kg, completa uma órbita do Sol a cada 1,2 anos, e se aproxima da Terra a cada 6 anos. Sua órbita é relativamente bem conhecida, e constantemente está sendo atualizada.

\f{caps01}{spaceflight101.com}{Esquema da cápsula de retorno de amostras}
\f{caps02}{spaceflight101.com}{Cápsula de retorno de amostra, durante testes}
\f{osir01}{spaceflight101.com}{Espaçonave e cápsula de retorno de amostra}

Para o retorno da amostra será utilizada uma cápsula (\rf{caps01}) baseada na que foi utilizada na missão Stardust da Nasa. A \rf{caps02} mostra a cápsula aberta durante testes, e a \rf{osir01} mostra a cápsula e o deck de instrumentos. Ela é uma peça crítica para o sucesso da missão, e tem a tarefa de transportar o material do asteroide até a Terra, reentrar na atmosfera protegendo contra o calor, e usar um paraquedas de aterrissagem.

A \rf{cron} mostra o cronograma do projeto OSIRIS-Rex com os principais eventos, desde o início em 2010, passando pelo retorno da amostra, até a fase final planejada do projeto que é a análise contínua de amostras.

\f{cron}{Nasa}{Crononograma do projeto OSIRIS-Rex}

O projeto da espaçonave é baseado nas missões MRO (Mars Reconnaissance Orbiter) e MAVEN (Mars Atmosphere e Volatile EvolutioN), aproveitando os principais componentes. Os recursos possuem margens de segurança, e os sistemas são redundantes. O sistema de voo inclui a estrutura e todos os vários componentes do subsistema para controlar e operar o veículo, o TAGSAM (Touch-And-Go Sample Acquisition Mechanism), o SRC ( Sample Return Capsule) e os cinco instrumentos científicos.

O EPS (Subsistema de Energia Elétrica) inclui dois painéis solares rígidos, manipulados sobre os eixos Y e Z da espaçonave. Além disso, duas baterias são utilizadas para manobras quando o sol não está incidindo nos painéis solares.

O subsistema de propulsão (PPS)  é monopropelente e tolerante a falhas, e inclui motores principais, propulsores de manobra de correção de trajetória, propulsores de sistema de controle de atitude e baixo conjuntos de motores de reação de empuxo. A propulsão está envolvida em todas as fases da missão, incluindo a fase de partida da Terra para ajustar a velocidade de escape da Terra; a fase de cruzeiro para ajustar a trajetória e garantir uma trajetória perfeitamente precisa, e a chegada em Bennu.

O subsistema GN\&C (Orientação, Navegação e Controle) inclui quatro RWAs (Reaction Wheel Assemblies) para realizar o giro da espaçonave e apontar os instrumentos durante as operações científicas. Essas rodas de reação também armazenam o momento do sistema entre os eventos. O subsistema GN\&C é responsável por comandar todos os propulsores da espaçonave, incluindo a execução de manobras de correção de trajetória. O subsistema utiliza uma IMU (Unidade de Medição Inercial) e rastreadores de estrelas para auxiliar o comando de atitude a bordo. Os sensores solares também suportam operações de segurança autônomas. Dois sensores GN\&C fornecem medições usadas para navegação.

O subsistema de comunicações utiliza banda X usando uma antena de alto ganho MAVEN e um amplificador MRO para \textit{downlink} com alta taxa de dados. Uma antena de ganho médio é utilizada durante a fase da missão TAG (``Touch-And-Go''). Além disso, duas antenas de baixo ganho estão disponíveis para TAG, e também são usadas para comando de \textit{downlink} e \textit{uplink} de dados de engenharia nominais
\cite{osiris2020}.


\section{Funcionamento}

\f{orbo}{Europa Press}{Órbita e trajetória da Osiris em torno do asteroide}

Em dezembro de 2018, após uma viagem de aproximadamente dois anos, a espaçonave encontrou o asteroide (\rf{orbo}), entrou em órbita, e durante meses recolheu informações e analisou a superfície, a uma distância de aproximadamente 5 km, para que a equipe da missão pudesse selecionar um local apropriado para extrair amostras. Quatro locais candidatos foram escolhidos e um selecionado. Em dezembro de 2020 a Osiris iniciou as manobras de aproximação, e coletou amostras. A amostra retornará à Terra em uma cápsula com paraquedas, e a viajem de volta será mais curta do que a de ida.

Antes da coleta da amostra, dois ensaios foram realizados, um em abril e outro em agosto, o primeiro chegou a 65 m da superfície, e o segundo a 40 m. Para a coleta, os painéis solares foram colocados em uma posição de 45\g em ``Y'' para diminuir a chance de acúmulo de poeira e também para oferecer mais distância ao solo evitando um possível contato. A aproximação foi lenta para minimizar os disparos do propulsor e evitar contaminação da amostra pelo propelente. Um acelerômetro indicou o contato com a superfície e uma mola no braço reduziu a força de impacto. Após o contato, um braço robótico e uma explosão usando nitrogênio foram utilizados para colher as amostras, e fotos foram tiradas para verificar se amostras foram recolhidas. Estava nos planos a possibilidade de uma nova tentativa, mas não foi necessário. Após a aquisição da amostra, o plano era  girar a espaçonave em torno do braço robótico para medir o momento de inércia e determinar a massa recolhida que deveria estar acima de 60 gramas, porém não foi realizado para evitar perda de amostra. O material da amostra foi então colocado na cápsula de retorno à Terra.

Assim que a amostra foi colocada na cápsula de retorno, foi feita uma \textit{verificação de retorno} utilizando o braço robótico para verificar se o coletor de amostra estava bem preso na cápsula de retorno. Após a verificação o braço foi solto do coletor, e a cápsula foi fechada e selada.

A janela de partida para o retorno à Terra será em março de 2021, e enquanto isso a espaçonave está sendo preparada. A cápsula de retorno está programada para chegar em setembro de 2023, fazer a reentrada na atmosfera, e pousar usando paraquedas no deserto de Utah.

\section{Equipamentos} %o que vai nele, detalhar

\f{ori}{\citeonline{NMOR}}{Instrumentos da OSIRIS-REx}
\f{ocams}{Wikimedia}{Grupo de câmeras de imagem - OCAMS}
\f{ovirs}{Wikimedia}{Espectrômetro Visível e IR - OVIRS}
\f{otes}{Wikimedia}{Espectrômetro de Emissão Termal - OTES}
\f{rexis}{Nasa}{Espectrômetro de imagens de raios X - REXIS}
\f{ola}{Nasa}{Espectrômetro de imagens de raios X - REXIS}
\f{tagsam}{Wikimedia}{Sistema de retorno de amostra - TAGSAM}

A espaçonave tem um conjunto de instrumentos de geração de imagens para analisar o asteroide em muitos comprimentos de onda, além do equipamento de telecomunicação (\rf{ori}).

O grupo de câmeras OCAMS (\rf{ocams}) coleta informações do asteroide fornecendo mapeamento global, reconhecimento e caracterização do local que será extraída a amostra, imagens de alta resolução, e registra (fotografa) a aquisição da amostra. É composto de 3 instrumentos, PolyCam, MapCam, e SamCam. 

O PolyCam é um telescópio de 20 cm, responsável pelas imagens de luz visível durante a aproximação do asteroide, e imagens de alta resolução da superfície. O MapCam é responsável pela procura de satélites e plumas de liberação de gás, mapeando em quatro canais (azul, verde, vermelho, e infravermelho). Fornece o modelo da forma do asteroide, e imagens de alta resolução dos locais que são candidatos a amostragem. O SamCam faz o monitoramento contínuo da aquisição da amostra.

O espectrômetro de luz visível e IR OVIRS (\rf{ovirs}) procura por minerais e substâncias orgânicas na superfície do asteroide. Fornece dados espectrais com resolução de 20 m. Mapeia do azul ao infravermelho (400-4300 nm) com resolução espectral de 7,5-22 nm. Estes dados são usados em conjunto com os espectros OTES para orientar a seleção do local da amostra.

O espectrômetro de emissão térmica OTES (\rf{otes}) cria um mapa espectral de emissão térmica e fornece informações espectrais de locais de coleta de amostra, usando um canal de infravermelho térmico de 4-50 $\mu$m para mapear minerais e materiais orgânicos. Ele também é usado para medir a emissão térmica total do asteroide.

O espectrômetro de imagens de raios X REXIS (\rf{rexis}) é responsável pelo mapa de espectroscopia de raios X do asteroide com a quantidade dos elementos presentes. É desenvolvido pelo MIT e pela Harvard e envolve mais de 100 alunos no desenvolvimento. É baseado em hardware e herança de voo, desta forma reduzindo os riscos técnicos, de cronograma, e de custo. Ele é um telescópio de raio-X entre 0,3-7,5 keV que fotografa emissões produzidas pelas absorção de raios-X solares e elementos do vento solar no regolito do asteroide que levam a emissões locais de raios-X. Em novembro de 2019 pesquisadores envolvidos na missão descobriram, usando o REXIS, rajadas de raio-X de um buraco negro.

O altímetro a laser OLA (\rf{ola}) é responsável pelas informações topográficas de alta resolução. Ele cria mapas topográficos globais do asteroide, mapas dos locais de amostra, e analisa navegação e gravidade. Ele varre a superfície em intervalos específicos para mapear rapidamente toda superfície. Os dados coletados também servem para aprimorar os estudos gravitacionais do asteroide. Ele é constituído de um receptor e dos transmissores complementares que aumentam a resolução das informações coletadas. O transmissor a laser de alta energia é usado no mapeamento de 1 a 7,5 km de distância. O de baixa energia é usado de 0,5 a 1 km. Os pulsos de laser são direcionados para um espelho de varredura móvel, que é co-alinhado com o telescópio receptor, limitando os efeitos da radiação solar de fundo. Cada pulso fornece as informações: intervalo de destino, azimute, elevação, intensidade recebida, e um marcador de tempo. Foi construído no Canadá.

O sistema de retorno de amostra, denominado \textit{Touch-And-Go Sample Acquisition Mechanism} TAGSAM (\rf{tagsam}) é um braço articulado contendo uma cabeça coletora de amostra. Ele tem uma fonte de nitrogênio para disparar contra a superfície e obter as amostras. O braço é responsável por coletar a amostra e colocá-la na cápsula de retorno à Terra.


%% EOF